\begin{prob}
    An understanding of data structures and the time complexity of the various
    operations involving them is crucial when designing algorithms. Using
    the wrong data structure, or using the right data structure in the wrong
    way, can lead to an algorithm which is still correct, but significantly
    slower. We will see one example in this problem.

    Another way of implementing the \python{merge} function used in
    \python{mergesort} is shown below:

    \begin{minted}[autogobble]{python}
        def merge(left, right, arr):
            left.append(float('inf'))
            right.append(float('inf'))
            for ix in range(len(arr)):
                if left[0] < right[0]:
                    arr[ix] = left.pop(0)
                else:
                    arr[ix] = right.pop(0)
    \end{minted}

    This code does essentially the same thing as the \python{merge} we saw in class:
    it merges two sorted lists \python{left} and \python{right} into a list
    \python{arr} such that the result is in sorted order. However, instead of
    using index variables to keep track of the front of \python{left} and \python{right},
    this code uses \python{list.pop(0)} to remove the first element from
    the appropriate list.

    \begin{subprobset}
        \begin{subprob}
            Popping an element from the end of a Python list is quick and
            efficient, but popping the first element from a list requires that
            the remaining elements be shifted over to the left by one place. Shifting
            requires copying the elements, one by one, to their new location.
            What is the time complexity of \python{arr.pop(0)}, assuming that
            \python{arr} is a list of size $n$?

            \begin{soln}
                $\Theta(n)$. Shifting one element to the left takes constant time.
                There are $n-1$ elements to shift, though, for a total time
                of $\Theta(n-1) = \Theta(n)$.
            \end{soln}
        \end{subprob}

        \begin{subprob}
            What is the time complexity of the new version of
            \python{merge}, assuming that (when they are passed to the
            function) \python{left} and \python{right} are lists of size
            $n/2$, and \python{arr} is a list of size $n$?

            \begin{soln}
                First, we claim that the calls to \python{.pop(0)} will dominate
                the time complexity. Therefore, we calculate the total amount
                of time spent in these calls.

                Each iteration of the \python{for}-loop will result in one
                list element being popped, either from \python{left} or
                from \python{right}. The time it takes to perform this pop
                will depend on how many elements are in \python{left} or
                \python{right} at that moment. But this depends on which
                list was popped from in earlier iterations, which
                ultimately depends upon the contents of \python{list} and
                \python{right}.  In short, there is no general way to say
                how large \python{left} or \python{right} are on the
                $\alpha$th iteration.

                But that's OK. Instead, recognize that there will be exactly
                one call to \python{left.pop(0)} when \python{left} has
                $n/2 + 1$ elements (remember that we appended $\infty$ to
                the end). Because popping takes time linear in the number
                of elements, this will take time $c(n/2 + 1)$, where $c$ is
                some constant. Some time later -- perhaps on the next
                iteration, perhaps in five iterations -- we will call
                \python{left.pop(0)} again. This will now take roughly $c\cdot n/2$
                time, since there are now $n/2$ elements in the list. On
                the third call, whenever that will be, the pop will take
                roughly $c(n/2 - 1)$ time. And so on.

                In total, calls to \python{left.pop(0)} will take time:
                \[
                    c \left[\left(n/2 + 1\right) + n/2 + (n/2 - 1) + \ldots +  2 + 1 \right]
                \]
                The sum of the first $k$ integers is $k(k+1)/2$. Here we have the sum
                of the first $k = (n/2 + 1)$ integers. Therefore, the total is:
                \[
                    c\cdot \frac{(n/2 + 1)(n/2 + 2)}{2} = \Theta(n^2).
                \]

                We haven't yet considered the calls to \python{right.pop(0)},
                but they will just double the total time found above. And
                since $2 \cdot \Theta(n^2)$ is still $\Theta(n^2)$, we
                see that the overall time taken is $\Theta(n^2)$.
            \end{soln}
        \end{subprob}

        \begin{subprob}
            What is the time complexity of a \python{mergesort} which uses the
            above version of \python{merge} instead of the linear time version
            used in lecture?

            \begin{soln}
                Recall that the recurrence describing the original version of
                \python{mergesort} was $T(n) = 2 T(n/2) + \Theta(n)$. There,
                the $\Theta(n)$ came from the call to \python{merge}. In the
                new version of \python{mergesort}, the $\Theta(n)$ is replaced
                by $\Theta(n^2)$. Hence we solve the recurrence relation:
                \[
                    T(n) =
                    \begin{cases}
                        2 T(n/2) + n^2, & \text{$n > 1$},\\
                        1,              & \text{otherwise}
                    \end{cases}
                \]

                We solve this by unrolling. We have $T(n/2) = 2 T(n/4) +
                (n/2)^2 = 2 T(n/4) + n^2/4$. Making this substitution:
                \begin{align*}
                    T(n) 
                        &= 2 T(n/2) + n^2\\
                        &= 2 \left[ 2 T(n/4) + n^2/4 \right] + n^2\\
                        &= 4 \cdot T(n/4) + n^2/2 + n^2\\
                        &= 4 \cdot T(n/4) + \frac{3n^2}{2}\\
                    \intertext{Unrolling again, we have $T(n/4) = 2 T(n/8) + n^2/16$. Hence:}
                        &= 4 \cdot \left[
                                2 T(n/8) + n^2 / 16
                            \right]
                            + \frac{3n^2}{2}\\
                        &=
                            8 T(n/8) + n^2/4 + \frac{3n^2}{2}\\
                        &=
                        8 T(n/8) + \frac{7n^2}{4}
                \end{align*}
                To summarize, we have:

                \begin{tabular}{cc}
                    Unroll \# & Formula\\
                    1   & $T(n) = 2 T(n/2) + n^2$\\
                    2   & $T(n) = 4 T(n/4) + \frac{3n^2}{2}$\\
                    3   & $T(n) = 8 T(n/8) + \frac{7n^2}{4}$\\
                \end{tabular}

                We propose the following general formula. On the $k$th unroll:
                \[
                    T(n) = 2^k T(n/2^k) + \frac{2^{k}-1}{2^{k-1}}n^2
                \]

                This process will continue until we have $n/2^k = 1$, that is,
                $\log_2 n$ times. Using this value of $k$ in our general formula,
                we obtain:

                \begin{align*}
                    T(n) 
                        &=
                            2^{\log_2 n} T(n/2^{\log_2 n}) + \frac{2^{\log_2 n} - 1}{2^{(\log_2 n) - 1}} n^2\\
                        &=
                            n T(n/n) + \frac{n - 1}{n/2} n^2
                            \\
                        &= n \cdot 1 + \Theta(1) \cdot n^2
                            \\
                        &=
                            \Theta(n^2)
                \end{align*}

                Hence the time complexity of this new algorithm is $\Theta(n^2)$.
                This is much worse than the time complexity of a well-implemented
                \python{mergesort}, which is $\Theta(n \log n)$.
            \end{soln}
        \end{subprob}

    \end{subprobset}
\end{prob}
